\begin{table}[t]
\centering
\caption{Standardized Effect Sizes}
\label{tab:sde}
\begin{adjustbox}{max width=\textwidth}
\begin{tabular}{lcccccc}
\toprule
Outcome & $\hat{\beta}$ & SE & SD($Y$) & SDE & SE(SDE) & Classification \\
\midrule
\multicolumn{7}{l}{\textit{Panel A: Pooled}} \\
log(Days Away + 1) & -0.0007 & 0.0011 & 1.8456 & -0.0032 & 0.0053 & Null \\
High Severity & -0.0006 & 0.0003 & 0.4270 & -0.0129 & 0.0057 & Small negative \\
Days Away (levels) & -0.0852 & 0.0355 & 63.2618 & -0.0116 & 0.0048 & Small negative \\
log(Days Away + 1), quadratic at mean & 0.0040 & 0.0022 & 1.8456 & 0.0188 & 0.0102 & Small positive \\
\addlinespace
\multicolumn{7}{l}{\textit{Panel B: Heterogeneous (Coal vs.\ Metal/Nonmetal)}} \\
Coal mines: log(Days Away + 1) & -0.0032 & 0.0025 & 1.9673 & -0.0121 & 0.0095 & Small negative \\
Metal/Nonmetal mines: log(Days Away + 1) & 0.0017 & 0.0014 & 1.6307 & 0.0096 & 0.0079 & Small positive \\
\bottomrule
\end{tabular}
\end{adjustbox}
\begin{tablenotes}[flushleft]
\item \textit{Notes:} \textbf{Country:} United States. \textbf{Research question:} Does mine-specific worker tenure reduce workplace injury severity, conditional on an accident occurring, beyond what general mining experience provides? \textbf{Policy mechanism:} Federal Mine Safety and Health Act requires operators to report individual accident details including worker experience at three levels (total mining, mine-specific, job-specific), enabling decomposition of establishment-specific vs.\ general human capital returns to safety. \textbf{Outcome definition:} Days away from work as reported on MSHA Form 7000-1; high severity includes fatalities, permanent disabilities, and injuries with days away. \textbf{Treatment:} Continuous; mine-specific tenure measured in years at the current mine at the time of injury. \textbf{Data:} MSHA accident records, 2000--2025, individual injury level, 218,953 observations after singleton removal. \textbf{Method:} OLS with mine, year, occupation, and subunit fixed effects; standard errors clustered at the mine level. \textbf{Sample:} All reported injuries with non-missing values for the three experience dimensions. SDE $= \hat{\beta} \times \text{SD}(X) / \text{SD}(Y)$ where SD($X$) is the standard deviation of mine-specific tenure and SD($Y$) is the standard deviation of the outcome. Classification refers to magnitude, not statistical significance: Large ($|$SDE$| > 0.15$), Moderate ($0.05$--$0.15$), Small ($0.005$--$0.05$), Null ($< 0.005$).
\end{tablenotes}
\end{table}

